% Copyright 2018  李文威 (Wen-Wei Li).
% Permission is granted to copy, distribute and/or modify this
% document under the terms of the Creative Commons
% Attribution 4.0 International (CC BY 4.0)
% http://creativecommons.org/licenses/by/4.0/

% 目的：字体相关设置，呼叫相关宏包。
% 将由 AJbook.cls 引入。
\ProvidesFile{font-setup-open.tex}[2026/03/01 Fandol font setup]

% ------------------------------------------------------------
% 西文无衬线字体
% ------------------------------------------------------------

\setsansfont{texgyreheros}[
  Extension      = .otf,
  UprightFont    = *-regular,
  BoldFont       = *-bold,
  ItalicFont     = *-italic,
  BoldItalicFont = *-bolditalic
]

% ------------------------------------------------------------
% 中文字体
% ------------------------------------------------------------

% 让 ctex 负责设置中文主字体。
% 下面不要再使用 \setCJKmainfont，否则会产生
% Redefining CJKfamily `\CJKrmdefault' 警告。
\usepackage[fontset=fandol]{ctex}

% 西里尔字符（如 Чеботарёв、Боревич）交给 CJK 字体渲染
\xeCJKDeclareCharClass{CJK}{"0400 -> "04FF}

% ------------------------------------------------------------
% 中文断行
% ------------------------------------------------------------

\XeTeXlinebreaklocale "zh"
\XeTeXlinebreakskip = 0pt plus 1pt minus 0.1pt

% ------------------------------------------------------------
% AJbook 需要的字体切换命令
% ------------------------------------------------------------

\newcommand{\thmheiti}{\heiti}
\newcommand{\sectionfont}{\bfseries\heiti}
\newcommand{\pffont}{\heiti}

% 保持 AJbook 原来的强调方式
\renewcommand{\em}{\bfseries\heiti}

% ------------------------------------------------------------
% AJbook 使用的额外 CJK 字体族
% ------------------------------------------------------------

% 日志显示 AJbook 中直接调用了：
%
%   \CJKfamily{song}
%
% 因此必须显式定义 song。
% 中文“斜体”映射到楷体，粗体映射到粗宋。
\setCJKfamilyfont{song}[
  BoldFont       = FandolSong-Bold,
  ItalicFont     = FandolKai-Regular,
  BoldItalicFont = FandolKai-Regular
]{FandolSong-Regular}

% 黑体别名
\setCJKfamilyfont{hei2}[
  BoldFont       = FandolHei-Bold,
  ItalicFont     = FandolHei-Regular,
  BoldItalicFont = FandolHei-Bold
]{FandolHei-Regular}

% 章节标题字体族
\setCJKfamilyfont{sectionfont}[
  BoldFont       = FandolHei-Bold,
  ItalicFont     = FandolHei-Regular,
  BoldItalicFont = FandolHei-Bold
]{FandolHei-Regular}

% 证明字体族
\setCJKfamilyfont{pffont}[
  BoldFont       = FandolHei-Bold,
  ItalicFont     = FandolHei-Regular,
  BoldItalicFont = FandolHei-Bold
]{FandolHei-Regular}

% 强调字体族
\setCJKfamilyfont{emfont}[
  BoldFont       = FandolHei-Bold,
  ItalicFont     = FandolHei-Regular,
  BoldItalicFont = FandolHei-Bold
]{FandolHei-Regular}
% ------------------------------------------------------------
% 让 \textit 中的西文使用斜体，中文使用楷体
% ------------------------------------------------------------

% 定义专供 \textit 使用的中文楷体族。
% 明确提供 m/it、b/n、b/it 映射，避免字体形状缺失警告。
\setCJKfamilyfont{textitkai}[
  BoldFont       = FandolKai-Regular,
  ItalicFont     = FandolKai-Regular,
  BoldItalicFont = FandolKai-Regular
]{FandolKai-Regular}

% 保存 LaTeX 原有的 \textit
\NewCommandCopy{\AJoriginaltextit}{\textit}

% 西文仍由原 \textit 排成斜体；
% 中文切换到上面定义的楷体族。
\RenewDocumentCommand{\textit}{m}{%
  \AJoriginaltextit{\CJKfamily{textitkai}#1}%
}

\endinput
